\chapter{Training Infrastructure}\label{ch:ml:training-infrastructure}

A team's order-book model trains for a day, and the team prices a second \omterm{def:ml:training-infrastructure:accelerator}{accelerator}. Before it buys one it should know
where the time goes. On the chapter's small version of the problem, measured on a laptop, a third of every training step is
spent assembling the batch, not computing on it; changing the storage format barely matters, and gathering the batch in one
indexed read instead of window by window cuts the data's share to a fifth. This chapter is about the machinery under the
models of Parts I to IV: \omterm{def:ml:training-infrastructure:accelerator}{accelerators} and where training time goes, \omterm{def:ml:training-infrastructure:loader}{data loaders} for tick data that never let a window cross
a session, lower-precision arithmetic, training on several devices, and checkpoints that let a killed run resume exactly.
Everything runs on one CPU thread, with no worker processes; the multi-device parts are simulated in one process, which is
enough to check their arithmetic, and the timings are the author's laptop's, labelled as such.

\section{Accelerators and where the time goes}

\begin{definition}[Accelerator]\label{def:ml:training-infrastructure:accelerator}
An \emph{accelerator}\index{accelerator} is a processor built for the dense linear algebra of \omterm{def:ml:neural-networks-for-noisy-tabular-data:mlp}{neural networks} (a graphics
processor or a dedicated tensor processor): thousands of arithmetic units, high-bandwidth memory beside them, and
low-precision matrix units that do most of the work.
\end{definition}

\begin{dated}{2026-09}{A data-centre accelerator}\label{dat:ml:training-infrastructure:h100}
NVIDIA's specification page for the H100 lists 80~GB of memory with 3.35~TB/s of bandwidth for the SXM version (94~GB and
3.9~TB/s for the NVL version), and 1\,979 teraflops of \omterm{def:ml:training-infrastructure:mixed}{bfloat16} tensor-core arithmetic for the SXM version, a figure the page
marks as with sparsity.
\end{dated}

Figures like these explain the usual bottleneck. An \omterm{def:ml:training-infrastructure:accelerator}{accelerator} can consume data far faster than a storage system, a
decompressor or a Python loop can feed it, so the first question about a slow training job is what fraction of each step the
device waits. The step divides into loading (reading the data from storage), batching (cutting windows and stacking them) and
computing (forward pass, backward pass, optimiser update). The chapter measures the three on its own job: eight ten-minute
\texttt{firm.tape} sessions of the book's synthetic market, each a table of 1\,140 half-second rows and 41 columns (the mid
price and the ten-level book of chapter~8), and a two-layer network that reads the last twenty rows and predicts the mid's
change over the next ten.

\section{Data loaders for tick data}

\begin{definition}[Data loader, memory-mapped dataset]\label{def:ml:training-infrastructure:loader}
A \emph{data loader}\index{data loader} is the code that turns stored data into training batches: it reads, decodes, samples,
cuts and stacks examples, ideally while the device computes on the previous batch. A \emph{memory-mapped
dataset}\index{memory-mapped dataset} stores the data in one binary file that the operating system maps into the process's
address space: nothing is read until a slice is touched, the operating system caches what is read, and several processes can
share one copy.
\end{definition}

Tick data impose one rule on the loader: a training window must not cross a session boundary, or a window will contain the
close of one day and the open of the next as if they were consecutive (and, with overlapping sessions for several instruments,
leak across them). The sampler of \cref{lst:ml:train:sampler} draws (session, start) pairs uniformly among the windows that fit,
and its random state is part of what a checkpoint saves.

The sessions are stored three ways: one CSV file per session (2.21~MB in all), one NumPy file per session (1.50~MB), and one
memory-mapped float32 file with a small JSON index (1.50~MB). \Cref{fig:ml:train:time} shows the measured cost of a training step
with each, loading being charged once per \omterm{def:ml:neural-networks-for-noisy-tabular-data:backprop}{epoch} of 100 steps. Parsing the CSV files takes 12.5~ms per \omterm{def:ml:neural-networks-for-noisy-tabular-data:backprop}{epoch} (176~MB/s),
reading the NumPy files 0.83~ms (1.8~GB/s), and mapping the single file 0.39~ms whatever its size. Per step these are tiny
against the half millisecond of cutting 256 windows one by one in Python and the 1.1 to 1.2~ms of computing: the data's share of
the step is 0.34, 0.31 and 0.39. Gathering the whole batch from the mapped file in one indexed read
(\cref{lst:ml:train:gather}) takes 0.27~ms and brings the share to 0.20.

\begin{omfigure}
\begin{tikzpicture}
\begin{axis}[width=11cm,height=6cm,ybar stacked,bar width=22pt,ymin=0,ymax=2.6,ylabel={milliseconds per training step},
  area legend,xtick={0,1,2,3},xticklabels={{CSV},{NumPy},{memory map},{memory map,\\one gather}},xticklabel style={align=center},
  tick label style={font=\small},label style={font=\small},xmin=-0.6,xmax=3.6,grid=major,grid style={gray!25},
  legend columns=3,legend style={font=\footnotesize,at={(0.5,-0.3)},anchor=north,draw=none,fill=none},table/col sep=comma]
\addplot[fill=omThm,draw=none] table[x=i,y expr=\thisrow{load_ms}/100]{figdata/ml/23-training-infrastructure/measured_loader.csv};
\addplot[fill=omProp,draw=none] table[x=i,y=batch_ms]{figdata/ml/23-training-infrastructure/measured_loader.csv};
\addplot[fill=omDef,draw=none] table[x=i,y=compute_ms]{figdata/ml/23-training-infrastructure/measured_loader.csv};
\legend{loading (per step),batching,computing}
\end{axis}
\end{tikzpicture}

\omcaption{Where a training step's time goes, by storage format and batch assembly: medians of five \omterm{def:ml:neural-networks-for-noisy-tabular-data:backprop}{epochs} of 100 steps of 256
windows. Measured on a laptop (Intel Core Ultra 7 155H) under WSL2, one thread, no isolated cores; other jobs may have run.
Data: \texttt{bench\_train.py}, \texttt{measured\_loader.csv}.}\label{fig:ml:train:time}
\end{omfigure}

At this scale the lesson is about code, not storage: the Python loop over windows costs more than any format. At the scale of a
year of order-book data the formats diverge. A year of the same half-second table for one instrument, 250 sessions of six and a
half hours, is about 1.9~GB of float32 and 2.8~GB of text: at the measured rates, parsing the text costs about 16~seconds per
\omterm{def:ml:neural-networks-for-noisy-tabular-data:backprop}{epoch}, reading the binary files about one second, and mapping costs nothing until the windows are read, from the page cache
after the first \omterm{def:ml:neural-networks-for-noisy-tabular-data:backprop}{epoch}. And beyond the memory of the machine only a memory-mapped or streamed format works at all.

\section{Mixed precision}

\begin{definition}[Mixed-precision training, bfloat16]\label{def:ml:training-infrastructure:mixed}
\emph{Mixed-precision training}\index{mixed-precision training} runs the expensive operations (matrix products, convolutions)
in a 16-bit format while keeping the weights, the loss and the optimiser's state in 32 bits (Micikevicius and co-authors,
2018). \emph{bfloat16}\index{bfloat16} is a 16-bit floating-point format with float32's eight exponent bits and seven fraction
bits: the same range as float32, and a machine epsilon (Book~4, chapter~25) of $2^{-7}\approx0.0078$ instead of $2^{-23}$, so no
loss scaling is needed against underflow (Kalamkar and co-authors, 2019).
\end{definition}

Under the CPU's \omterm{def:ml:training-infrastructure:mixed}{bfloat16} autocast the chapter's network trains to the same place (\cref{fig:ml:train:bf16}): over the last 50
of 300 steps the mean loss is 0.0525 against 0.0527 in float32. Step by step the two runs differ: the largest difference between
their losses over the run is 0.0132, 22\% of that step's float32 loss, because rounding changes each update slightly and the
trajectories drift apart. Reproducing a run bit for bit across precisions is not to be expected; equivalence is judged on the
validation metric, over the run.

\begin{omfigure}
\begin{tikzpicture}
\begin{axis}[width=11cm,height=5.6cm,xlabel={training step},ylabel={loss (10-step average)},tick label style={font=\small},
  label style={font=\small},xmin=10,xmax=300,grid=major,grid style={gray!25},legend cell align=left,
  legend style={font=\footnotesize,at={(0.98,0.98)},anchor=north east,fill=white,draw=none},table/col sep=comma]
\addplot[omDef,very thick] table[x=step,y=fp32]{figdata/ml/23-training-infrastructure/losses.csv};
\addplot[omThm,thick,dashed] table[x=step,y=bf16]{figdata/ml/23-training-infrastructure/losses.csv};
\legend{float32,bfloat16 autocast}
\end{axis}
\end{tikzpicture}

\omcaption{Training loss of the same network, same data order and same initialisation, in float32 and under \omterm{def:ml:training-infrastructure:mixed}{bfloat16} autocast.
Data: \texttt{ml\_train.run} (\texttt{fig\_train.py}).}\label{fig:ml:train:bf16}
\end{omfigure}

\section{Distributed training}

\begin{definition}[Gradient accumulation, data parallelism, model parallelism, all-reduce]\label{def:ml:training-infrastructure:dist}
\emph{Gradient accumulation}\index{gradient accumulation} sums the gradients of several micro-batches before one optimiser step,
which equals one step on their union. \emph{Data parallelism}\index{data parallelism} runs a copy of the model on each device,
each on its own share of the batch, and averages their gradients before every step (Li and co-authors, 2020); \emph{model
parallelism}\index{model parallelism} splits one model's layers or matrices across devices when it does not fit on one
(Shoeybi and co-authors, 2019). An \emph{all-reduce}\index{all-reduce} is the collective operation that leaves every device with
the sum (or average) of all devices' arrays.
\end{definition}

Both identities are checked in the chapter's code (\cref{lst:ml:train:loop}). Four accumulated micro-batches of 64 and one batch
of 256 of the same windows leave parameters that differ by at most $1.4\times10^{-7}$ after 20 steps, rounding in a different
order of summation. The average of two workers' gradients on the halves of a batch differs from the whole batch's gradient by at
most $1.5\times10^{-8}$, against gradients of size 0.12: an \omterm{def:ml:training-infrastructure:dist}{all-reduce} computes exactly the large-batch gradient, and \omterm{def:ml:training-infrastructure:dist}{data
parallelism} is \omterm{def:ml:training-infrastructure:dist}{gradient accumulation} across devices. What changes with a larger effective batch is the optimisation: the \omterm{def:ml:trees-and-boosting:boosting}{learning
rate} and its warm-up must be rescaled (Goyal and co-authors, 2017), and for the small, noisy models of trading, larger batches
rarely help.

\section{Checkpoints and deterministic resumption}

\begin{definition}[Training checkpoint]\label{def:ml:training-infrastructure:ckpt}
A \emph{training checkpoint}\index{training checkpoint} is a saved state from which a run can continue as if it had not stopped:
the model's parameters, the optimiser's state (moment estimates, step counts), the data sampler's position and random state,
the random generators' states and the step count.
\end{definition}

A run of 200 steps and a run killed at step 100, checkpointed (\cref{lst:ml:train:ckpt}), and resumed into newly constructed
objects with a different initialisation end with bitwise identical parameters, and identical losses step for step; the checkpoint
takes 693~kB. Saving only the model's weights, the usual shortcut, restarts \omterm{def:ml:neural-networks-for-noisy-tabular-data:adam}{Adam}'s moment estimates and replays the sampler's
first windows; the resumed run ends with parameters up to 0.039 away from the uninterrupted one's. For research that must be
reproduced (chapter~25) and for jobs on machines that can be pre-empted, the full checkpoint is the difference between resuming
and restarting with a different run.

\begin{method}[Before buying an accelerator]\label{met:ml:train:method}
\begin{enumerate}
\item Measure load, batch and compute time per step; if data take more than a small share, fix the loader first (binary or
memory-mapped storage, batch gathering in one read, prefetching).
\item Try \omterm{def:ml:training-infrastructure:mixed}{bfloat16} autocast and judge it on the validation metric over the run.
\item Grow the batch by accumulation before growing the hardware; rescale the \omterm{def:ml:trees-and-boosting:boosting}{learning rate} and check that the model improves.
\item Checkpoint everything that makes the run's future, and test resumption bitwise.
\end{enumerate}
\end{method}

\section{Tutorial: where the time goes}\label{tut:ml:training-infrastructure:tut}

\begin{tutorial}
\textbf{Goal.} Store tick sessions three ways, measure a step's time budget, train in \omterm{def:ml:training-infrastructure:mixed}{bfloat16}, verify accumulation and \omterm{def:ml:training-infrastructure:dist}{data
parallelism}, and resume a run bitwise. \textbf{End state:} \cref{fig:ml:train:time,fig:ml:train:bf16}.
\begin{tutsteps}
\item \textbf{A session-aware window sampler with a saveable state.}
\omcode{code/firm/mltrain/firm_mltrain.py}{78}{96}{Windows that never cross a session boundary.}\label{lst:ml:train:sampler}
\item \textbf{One gather for a whole batch.}
\omcode{code/firm/mltrain/firm_mltrain.py}{110}{116}{A batch from the memory-mapped store in one indexed read.}\label{lst:ml:train:gather}
\item \textbf{Accumulation and \omterm{def:ml:training-infrastructure:mixed}{bfloat16} autocast.}
\omcode{code/firm/mltrain/firm_mltrain.py}{119}{141}{The training loop.}\label{lst:ml:train:loop}
\item \textbf{Checkpoints.}
\omcode{code/firm/mltrain/firm_mltrain.py}{156}{167}{Saving and restoring everything a run's future depends on.}\label{lst:ml:train:ckpt}
\item \textbf{Run} \texttt{ml\_train.precision()}, \texttt{accumulation()}, \texttt{data\_parallel()}, \texttt{resume()},
\texttt{naive\_resume()}, \texttt{fig\_train.py}; and, once, \texttt{bench\_train.py} under \texttt{nice}.
\end{tutsteps}
\textbf{What to change next.} Prefetch the next batch in a background thread and measure the overlap; add a second instrument
and check that windows never mix instruments.
\end{tutorial}

\section{Build: the training toolkit}\label{bld:ml:training-infrastructure:mltrain}

\begin{build}
\bfield{Purpose} Training runs that are fast where it matters, correct at session boundaries, and exactly resumable.
\bfield{Interface} \texttt{write\_shards(sessions, root, fmt)}, \texttt{Shards(root, fmt)}, \texttt{WindowSampler(lengths, T,
horizon, batch, seed)} with \texttt{state}/\texttt{set\_state}, \texttt{make\_batch}, \texttt{make\_batch\_vec},
\texttt{train(model, shards, sampler, steps, lr, accumulate, bf16)}, \texttt{data\_parallel\_grads}, \texttt{save\_checkpoint},
\texttt{load\_checkpoint}.
\bfield{Rules} Windows never cross sessions; timings are measured, labelled with the machine, and never asserted; everything that
determines the rest of a run goes into its checkpoint.
\bfield{Acceptance tests} \texttt{code/firm/mltrain/tests/}: the three formats return identical windows; no sampled window crosses
a boundary; the gathered batch equals the looped one; accumulation equals a larger batch; the averaged shard gradients equal the
full gradient; resumption is bitwise.
\bfield{Stretch} Real two-process \omterm{def:ml:training-infrastructure:dist}{data parallelism} with the gloo backend when the machine allows; streaming from compressed
shards; prefetching.
\end{build}

\begin{omsources}
\item P.~Micikevicius and co-authors, ``Mixed precision training'', arXiv:1710.03740, 2018.
\item D.~Kalamkar and co-authors, ``A study of BFLOAT16 for deep learning training'', arXiv:1905.12322, 2019.
\item S.~Li and co-authors, ``PyTorch distributed: experiences on accelerating data parallel training'', arXiv:2006.15704, 2020.
\item P.~Goyal and co-authors, ``Accurate, large minibatch SGD: training ImageNet in 1 hour'', arXiv:1706.02677, 2017.
\item M.~Shoeybi and co-authors, ``Megatron-LM: training multi-billion parameter language models using model parallelism'',
arXiv:1909.08053, 2019.
\item NVIDIA, H100 Tensor Core GPU specifications (web page, accessed September 2026).
\end{omsources}

\section{Exercises}

\begin{exercise}[$\star$]\label{exo:ml:training-infrastructure:1}
What is the machine epsilon of \omterm{def:ml:training-infrastructure:mixed}{bfloat16}, float16 and float32? Which of them can represent $10^{-30}$?
\end{exercise}

\begin{exercise}[$\star$]\label{exo:ml:training-infrastructure:2}
A session has 1\,140 rows; windows are 20 rows with a 10-row target horizon. How many windows does it hold? What goes wrong if a
sampler draws starts uniformly over the concatenated sessions?
\end{exercise}

\begin{exercise}[$\star$]\label{exo:ml:training-infrastructure:3}
From \cref{fig:ml:train:time}, what speed-up would halving the compute time give with the looped loader, and with the gather?
\end{exercise}

\begin{exercise}[$\star\star$]\label{exo:ml:training-infrastructure:4}
Why do float32 and \omterm{def:ml:training-infrastructure:mixed}{bfloat16} losses differ by up to 22\% at some steps while their averages agree?
\end{exercise}

\begin{exercise}[$\star\star$]\label{exo:ml:training-infrastructure:5}
Show that with a mean loss, accumulating $k$ micro-batches each scaled by $1/k$ gives the gradient of the mean over their union.
What changes with \omterm{def:ml:neural-networks-for-noisy-tabular-data:dropout}{batch normalisation}?
\end{exercise}

\begin{exercise}[$\star\star$]\label{exo:ml:training-infrastructure:6}
\emph{Find the flaw.} ``We resume pre-empted jobs from the last saved weights, so our runs are reproducible.''
\end{exercise}

\begin{exercise}[$\star\star\star$]\label{exo:ml:training-infrastructure:7}
Estimate, from the measured throughputs, the \omterm{def:ml:neural-networks-for-noisy-tabular-data:backprop}{per-epoch} loading time of a year of the chapter's table in each format (250
sessions of six and a half hours of half-second rows, 41 float32 columns, text about 1.5 times larger than binary).
\end{exercise}

\begin{exercise}[$\star\star\star$]\label{exo:ml:training-infrastructure:8}
In a ring \omterm{def:ml:training-infrastructure:dist}{all-reduce} over $p$ devices each holding $n$ numbers, each device sends and receives about $2n(p-1)/p$ numbers. Explain
why, and what it implies for how data-parallel training scales with the number of devices.
\end{exercise}

\section{Problem: Where the Time Goes}

\begin{problem}[{Weekend problem --- measure before you buy}]\label{pb:ml:training-infrastructure:1}
The chapter's sessions, formats, network and runs.

\medskip\noindent\textbf{Part I --- The budget.}
\begin{enumerate}
\item How large is the data in each format?
\item What are the load, batch and compute times per step?
\item What share of the step goes to data for each storage format, and with the gather?
\item What would you change first, and why?
\end{enumerate}

\medskip\noindent\textbf{Part II --- Precision.}
\begin{enumerate}[resume]
\item What are \omterm{def:ml:training-infrastructure:mixed}{bfloat16}'s bits, range and epsilon?
\item How do the float32 and \omterm{def:ml:training-infrastructure:mixed}{bfloat16} runs compare?
\item Why is no loss scaling needed with \omterm{def:ml:training-infrastructure:mixed}{bfloat16}?
\item Where would you keep float32 in a trading model's pipeline?
\end{enumerate}

\medskip\noindent\textbf{Part III --- Parallelism and checkpoints.}
\begin{enumerate}[resume]
\item What do the accumulation and data-parallel checks show?
\item What must a checkpoint contain, and what does leaving part of it out cost?
\item How large is the chapter's checkpoint?
\item What does a larger effective batch require of the \omterm{def:ml:trees-and-boosting:boosting}{learning rate}?
\end{enumerate}

\medskip\noindent\textbf{Part IV --- The verdict.}
\begin{enumerate}[resume]
\item State the \emph{named result}: the data-loading share of step time for each storage format, and the largest loss difference
between float32 and \omterm{def:ml:training-infrastructure:mixed}{bfloat16} training over the run.
\item Would the team need a second \omterm{def:ml:training-infrastructure:accelerator}{accelerator}?
\item When does \omterm{def:ml:training-infrastructure:dist}{model parallelism} become necessary for a trading model?
\item How would you make a training run reproducible across machines?
\item What is the risk of a sampler that crosses sessions?
\item How would you measure a loader on shared infrastructure?
\item What would you log about every training run?
\item In one sentence: what is the first thing to measure about a slow training job?
\end{enumerate}
\end{problem}

\section{Interview questions}

\begin{interviewq}[$\star$ \iqroles{mle}]\label{iq:ml:training-infrastructure:1}
What is \omterm{def:ml:training-infrastructure:mixed}{bfloat16} and why is it preferred to float16 for training?
\end{interviewq}

\begin{interviewq}[$\star\star$ \iqroles{mle, developer}]\label{iq:ml:training-infrastructure:2}
Your GPU utilisation is 30\%. How do you find out why?
\end{interviewq}

\begin{interviewq}[$\star\star$ \iqroles{mle}]\label{iq:ml:training-infrastructure:3}
Explain \omterm{def:ml:training-infrastructure:dist}{data parallelism} and what an \omterm{def:ml:training-infrastructure:dist}{all-reduce} does.
\end{interviewq}

\begin{interviewq}[$\star\star$ \iqroles{mle, researcher}]\label{iq:ml:training-infrastructure:4}
How would you store and sample a year of order-book data for training sequence models?
\end{interviewq}

\begin{interviewq}[$\star\star$ \iqroles{mle}]\label{iq:ml:training-infrastructure:5}
What does a \omterm{def:ml:training-infrastructure:ckpt}{training checkpoint} need to contain for bitwise resumption?
\end{interviewq}

\begin{interviewq}[$\star\star\star$ \iqroles{mle, developer}]\label{iq:ml:training-infrastructure:6}
\omterm{def:ml:training-infrastructure:dist}{Gradient accumulation} and \omterm{def:ml:training-infrastructure:dist}{data parallelism} are said to be equivalent. When are they not?
\end{interviewq}
